\section*{Chapter \ref{ch:dv:options-market-making-in-practice} --- Options Market Making in Practice}

\begin{solution}{exo:dv:options-market-making-in-practice:1}
With no informed flow the profit is $20\,e^{-w/0.3}w$. Its derivative $20\,e^{-w/0.3}(1-w/0.3)$ vanishes at $w=0.3$, where the gain per trade and the loss of
trades balance at the margin.
\end{solution}

\begin{solution}{exo:dv:options-market-making-in-practice:2}
$0.42\times0.115=0.048$ and $0.42\times0.199=0.084$: 4.8 and 8.4 cents a share, \$4.83 and \$8.37 a contract.
\end{solution}

\begin{solution}{exo:dv:options-market-making-in-practice:3}
$2q=65\,894$, $fN=3\,000$, $N+2q=75\,894$: $65\,894\times3\,000/75\,894=2\,605$.
\end{solution}

\begin{solution}{exo:dv:options-market-making-in-practice:4}
The net is $g\,2q\,fN/(N+2q)-cq=aq/(N+2q)-cq$ with $a=2fNg$. Its derivative is $aN/(N+2q)^2-c$, zero at $(N+2q)^2=aN/c$, so
$q^*=(\sqrt{aN/c}-N)/2$; the second derivative is negative. With $a=288\,000$, $N=10\,000$, $c=0.5$: $q^*=(75\,895-10\,000)/2=32\,947$.
\end{solution}

\begin{solution}{exo:dv:options-market-making-in-practice:5}
Trading to the edge is the smallest trade that restores an acceptable risk. Trading to the centre pays for a larger trade and resets the position to where the
next random move is as likely to need another trade as not. Inside the band the risk is accepted by design, so paying to reduce it further is waste.
\end{solution}

\begin{solution}{exo:dv:options-market-making-in-practice:6}
The error is the skew times the log move since the fit, divided by $\sqrt\tau$. The log move's standard deviation grows as the square root of the interval, and
the skew's effect on a strike's volatility is larger for short expiries, where the smile is steeper in \omterm{def:dv:implied-volatility-and-its-surface:surface}{log-moneyness}: hence $1/\sqrt\tau$.
\end{solution}

\begin{solution}{exo:dv:options-market-making-in-practice:7}
The market makers' assigned calls are hypergeometric: $65\,894$ of $75\,894$ short positions, $72\,894$ exercises. With 200\,000 draws the net profit averages
\$108\,552 with a standard deviation of \$870 (\$871 from the hypergeometric formula); its first percentile is \$106\,502. Assignment risk is small at this size.
\end{solution}

\begin{solution}{exo:dv:options-market-making-in-practice:8}
One path. Over 20 seeds of the same toy the difference averages \$9\,585 with a standard deviation of \$11\,597 (a standard error of \$2\,593), and limits did worse
in 6 of the 20. Limits reduce inventory risk, and on average in this toy that also helps the P\&L, because the flow is one-sided. But a single month is a
single draw.
\end{solution}

\begin{solution}{pb:dv:options-market-making-in-practice:1}
\textbf{1.} Yes: the dividend (\$50 a contract) exceeds the put's value (\$2); failing costs \$48 a contract.
\textbf{2.} With zero rates a call's early exercise gains nothing except the dividend, which the holder of the share receives and the holder of the call does not.
\textbf{3.} More than half of the long positions that should have been exercised were not (1996--2006), costing holders over \$491 million.
\textbf{4.} Other short positions, including dealers'. Their shorts are assigned less than they would be if every holder exercised, so they keep the dividend
on the unassigned part.
\textbf{5.} The market makers add $2q$ short positions that share the assignments with the original writers. As $q$ grows, their shorts become almost all of
the short interest, so almost all the non-assignment falls to them.
\textbf{6.} $2q\,fN/(N+2q)$.
\textbf{7.} $fN=3\,000$ contracts, approached as $q\to\infty$.
\textbf{8.} \$0.50: four trades at \$0.10 and two exercises at \$0.05.
\textbf{9.} 32\,947 contracts each way; \$108\,553 net (2\,605 unassigned, \$125\,026 gross, \$16\,474 fees).
\textbf{10.} \$28\,591 (16\,909 contracts) and \$193\,510 (43\,990 contracts).
\textbf{11.} The gain shrinks; the standard deviation from random assignment is about \$870 at the best size.
\textbf{12.} Fewer failures: the gain falls, and at the planned size the fees can exceed it (at 10\% failure the best size is half as large).
\textbf{13.} Fees rise linearly with $q$, so the best size and the profit shrink; with no fees the play would capture all \$144\,000 at an unlimited size.
\textbf{14.} The share position of exercised calls and any imbalance of assignment, carried through the ex-date open, and operational risk on exercise instructions.
\textbf{15.} From the series' history of exercises on past ex-dates, the share of open interest held by retail accounts, and how deep in the money it is.
\textbf{16.} The holders who fail to exercise. The market makers do not cause the failure, but they redirect its value from the original writers to themselves.
\textbf{17.} Only at a size that pays the fees: at 10\% failure the best trade nets \$28\,591 on fees of about \$8\,500, still worth it on these numbers.
\textbf{18.} Automatic exercise of calls whose dividend exceeds the time value, or fees on dividend strategies high enough to make $q^*$ zero.
\textbf{19.} With 30\% of 10\,000 public contracts unexercised, a \$48 gain per unassigned call and \$0.50 of fees per contract traded each way: \$108\,553 at the best
size of 32\,947 contracts (\$28\,591 at 10\%, \$193\,510 at 50\%).
\textbf{20.} From option holders who fail to exercise, collected through the pro-rata assignment rule.
\end{solution}

\begin{solution}{iq:dv:options-market-making-in-practice:1}
In volatility points, then times \omterm{def:dv:greeks-and-the-hedging-pnl:greeks}{vega}: wide enough to cover the theo's error against informed traders, narrow enough to keep uninformed flow; estimated from the
mark-outs of past trades by counterparty type and adjusted for inventory and limits.
\iqlookfor{error of the theo, adverse selection, \omterm{def:dv:greeks-and-the-hedging-pnl:greeks}{vega} scaling, mark-outs.}
\end{solution}

\begin{solution}{iq:dv:options-market-making-in-practice:2}
Not on a clock: in a band around the target delta whose width grows with the cost of hedging and shrinks with gamma and risk aversion (half-width proportional
to the cube root of cost times gamma squared), trading to the band's edge. On the chapter's example a band gives the same cost as daily hedging with 41\% less
P\&L noise.
\iqlookfor{the band, and why it beats a clock.}
\end{solution}

\begin{solution}{iq:dv:options-market-making-in-practice:3}
On the eve of an ex-date, market makers trade large offsetting amounts of a deep in-the-money call and exercise their longs. Pro-rata assignment leaves most of the
public's unexercised calls assigned against someone else, so the market makers keep the dividend on their unassigned shorts. The holders who fail to exercise
lose, and the original writers lose what they would have kept.
\iqlookfor{pro-rata assignment, failure to exercise, who loses.}
\end{solution}

\begin{solution}{iq:dv:options-market-making-in-practice:4}
Fast: the forward and each series' theo and delta from the current surface parameters, by closed forms and cached per-expiry quantities, with the marking rule
between refits; quote updates when the theo moves by more than a threshold. Slow: surface refits (on a clock and on events), limit checks, width parameters from
mark-outs. The quotes are pulled while a refit runs.
\iqlookfor{separation of hot path and refit, event-driven refits.}
\end{solution}

\begin{solution}{iq:dv:options-market-making-in-practice:5}
Offsetting \omterm{def:dv:greeks-and-the-hedging-pnl:greeks}{vegas} in different expiries are exposed to changes in the term structure, which total \omterm{def:dv:greeks-and-the-hedging-pnl:greeks}{vega} hides. Near a limit the desk shades its quotes to attract
offsetting trades, hedges with listed options in the bucket, and withdraws the side that would breach the limit only at the limit.
\iqlookfor{term-structure risk, shading before withdrawing.}
\end{solution}

\begin{solution}{iq:dv:options-market-making-in-practice:6}
That is pin risk: uncertain assignment. Reduce the position before the close (buy back the strike or trade the pin), size the weekend share position to the expected
assignment, and keep capacity to hedge Monday's open; know the exercise cut-offs and the automatic-exercise threshold.
\iqlookfor{assignment uncertainty and acting before the close.}
\end{solution}
